% Ad Familiares 10.9 (ad-familiares-10-09) --- English translation
% Translated by Claude (Anthropic) from the Latin (Perseus).
% Style guide: STYLE.md.

\ciceroLetterOpener{PLANCUS CICERONI S.}

\ciceroSection{1} I am glad that you have undertaken nothing on my account rashly, and that you have undertaken nothing in vain when you pledged me to the rest. You have at any rate this much greater proof of my devotion: I wished my plans to be known to you sooner than to anyone else. By the day, I trust, you can see the additions to my services growing; that you will come to recognize them still more, I undertake.

\ciceroSection{2} As for myself, my dear Cicero --- so may the state be freed by my help from the evils now hanging over it! --- I look up to the honours and rewards you can bestow as something to be ranked, surely, with immortality, and yet without any one of them I shall not slacken in the slightest from my zeal and perseverance. Unless, in the great number of our best citizens, the drive of my spirit proves singular and my service outstanding, I want no addition to my standing from your votes.

\ciceroSection{3} As for ambition --- I covet nothing for myself (I fight against the thing); and I gladly leave you the regulator both of the timing and of the substance. Nothing granted to a citizen by his country can ever seem late or stinted. I crossed the Rhone with the army by forced marches on the sixth day before the Kalends of May; from Vienna I sent ahead a thousand horse by the shorter road. As for myself, if Lepidus does not obstruct me, I shall give satisfaction by my speed; if, however, he plants himself across my march, I shall take counsel as the moment requires. The forces I am bringing up are very strong in numbers, in quality, and in loyalty. I ask you to keep loving me, if you know I shall do the like in return. Farewell.
